\subsection*{Base de datos}
\addcontentsline{toc}{subsection}{Base de datos}

\us[
	identificador=US-BD-01,
	titulo={Gestor de migraciones},
	como={Desarrollador},
	quiero={Modificar el esquema de la base de datos de forma versionada junto al código que lo requiere, para que cualquier entorno pueda reconstruirse aplicando la misma secuencia de cambios},
	requerimientos={\reqref{RNF-DE-04}},
	puntos={2},
	criterios={
			\begin{itemize}
				\item El proyecto incluye un gestor de migraciones para la base de datos de la API.
				\item Las migraciones se almacenan en el repositorio, junto al código de la API.
				\item Existen comandos para crear una migración nueva y para aplicar las migraciones pendientes sobre la base de datos local.
				\item El gestor de migraciones obtiene los datos de acceso a la base de datos de las mismas variables de entorno que la API.
				\item Existe documentación sobre la creación y la aplicación de migraciones.
			\end{itemize}
		},
	tareas={
			\begin{itemize}
				\item Selección del gestor de migraciones.
				\item Instalación y configuración del gestor de migraciones en la API.
				\item Creación de los comandos para crear y aplicar migraciones en local.
				\item Actualización de la documentación del repositorio.
			\end{itemize}
		}
]
